\subsection{Groups Acting on Themselves by Left Multiplication---Cayley's Theorem}

\begin{exercise}[4.2.1]
	Let $G$ be $\set{1,a,b,c}$, the Klein 4-group whose group table is written out in Section 2.5.
	\begin{subexercises}
		\item Label $1,a,b,c$ with the integers $1,2,4,3$, respectively, and prove that under the left regular representation of $G$ into $S_4$ the non-identity elements are mapped as follows:
		\[
			a \mapsto (1\ 2)(3\ 4), \quad b \mapsto (1\ 4)(2\ 3), \quad c \mapsto (1\ 3)(2\ 4).
		\]
		\item Relabel $1,a,b,c$ as $1,4,2,3$, respectively, and compute the image of each element of $G$ under the left regular representation of $G$ into $S_4$.
		Show that the image of $G$ in $S_4$ under this labelling is the same \emph{subgroup} as the image of $G$ in part (a) (even though the non-identity elements individually map to different permutations under the two different labelling).
	\end{subexercises}
\end{exercise}

\begin{solenum}
	\item We have $a \cdot 1 = a$ so that $\sigma_a(1) = 2$.
	Similarly, $\sigma_a(2) = 1$.
	We also have $a \cdot b = c$ so that $\sigma_a(4) = 3$, and $\sigma_a(3) = 4$, and we have $a \mapsto (1\ 2)(3\ 4)$.
	The others are computed similarly. 
	\item Again, we have $a \cdot 1 = a$ so that $\sigma_a(1) = 4$.
	Similarly, $\sigma_a(4) = 1$.
	We also have $a \cdot b = c$ so that $\sigma_a(2) = 3$, and $\sigma_a(3) = 2$, and we have $a \mapsto (1\ 4)(2\ 3)$.
	For $b$, we have $b \mapsto (1\ 2)(3\ 4)$, and for $c$, we have $c \mapsto (1\ 3)(2\ 4)$.
	The image of $G$ in $S_4$ under this labelling is $\set{1, (1\ 4)(2\ 3), (1\ 2)(3\ 4), (1\ 3)(2\ 4)}$, which is the same subgroup as in part (a).
\end{solenum}

\begin{exercise}[4.2.2]
	List the elements of $S_3$ as $1, (1\ 2), (2\ 3), (1\ 3), (1\ 2\ 3), (1\ 3\ 2)$ and label these with the integers $1,2,3,4,5,6$ respectively.
	Exhibit the image of each element of $S_3$ under the left regular representation of $S_3$ into $S_6$.
\end{exercise}

\begin{sol}
	Consider $(1\ 2)$.
	Then we have the following computations:
	\begin{align*}
		(1\ 2) \cdot 1 & = (1\ 2) && \sigma_{(1\ 2)}(1) = 2 \\
		(1\ 2) \cdot (1\ 2) & = 1 && \sigma_{(1\ 2)}(2) = 1 \\
		(1\ 2) \cdot (2\ 3) & = (1\ 3\ 2) && \sigma_{(1\ 2)}(4) = 5 \\
		(1\ 2) \cdot (1\ 3) & = (1\ 2\ 3) && \sigma_{(1\ 2)}(3) = 6 \\
		(1\ 2) \cdot (1\ 2\ 3) & = (2\ 3) && \sigma_{(1\ 2)}(5) = 3 \\
		(1\ 2) \cdot (1\ 3\ 2) & = (1\ 3) && \sigma_{(1\ 2)}(6) = 4
	\end{align*}
	Thus, we have $(1\ 2) \mapsto (1\ 2)(3\ 6\ 4\ 5)$.
	We calculate that $(2\ 3) \mapsto (1\ 3\ 5\ 2)(4\ 6)$.
	Since $(1\ 2)(2\ 3) = (1\ 2\ 3)$, then
	\[
		(1\ 2\ 3) \mapsto (1\ 2)(3\ 6\ 4\ 5)(1\ 3\ 5\ 2)(4\ 6) = (1\ 5\ 6)(2\ 4\ 3)
	\]
	Continuing in this way, we find that the images of the elements of $S_3$ under the left regular representation are as follows:
	\begin{align*}
		1 & \mapsto 1 \\
		(1\ 2) & \mapsto (1\ 2)(3\ 6\ 4\ 5) \\
		(2\ 3) & \mapsto (1\ 3\ 5\ 2)(4\ 6) \\
		(1\ 3) & \mapsto (1\ 4)(2\ 6)(3\ 5) \\
		(1\ 2\ 3) & \mapsto (1\ 5\ 6)(2\ 4\ 3) \\
		(1\ 3\ 2) & \mapsto (1\ 6\ 5)(2\ 3\ 4) \qh
	\end{align*}
\end{sol}

\begin{exercise}[4.2.3]
	Let $r$ and $s$ be the usual generators for the dihedral group of order 8.
	\begin{subexercises}
		\item List the elements of $D_8$ as $1, r, r^2, r^3, s, sr, sr^2, sr^3$ and label these with the integers $1,2,\dots,8$ respectively.
		Exhibit the image of each element of $D_8$ under the left regular representation of $D_8$ into $S_8$.
		\item Relabel this same list of elements of $D_8$ with the integers $1,3,5,7,2,4,6,8$ respectively and recompute the image of each element of $D_8$ under the left regular representation with respect to this new labelling.
		Show that the two subgroups of $S_8$ obtained in parts (a) and (b) are different.
	\end{subexercises}
\end{exercise}

\begin{solenum}
	\item We calculate $\sigma_r = (1\ 2\ 3\ 4)(5\ 8\ 7\ 6)$ and $\sigma_s = (1\ 5)(2\ 6)(3\ 7)(4\ 8)$.
	Continuing in this way, we find that the images of the elements of $D_8$ under the left regular representation are as follows:
	\begin{align*}
		1 & \mapsto 1 & s & \mapsto (1\ 5)(2\ 6)(3\ 7)(4\ 8) \\
		r & \mapsto (1\ 2\ 3\ 4)(5\ 8\ 7\ 6) & sr & \mapsto (1\ 6)(2\ 7)(3\ 8)(4\ 5) \\
		r^2 & \mapsto (1\ 3)(2\ 4)(5\ 7)(6\ 8) & sr^2 & \mapsto (1\ 7)(2\ 8)(3\ 5)(4\ 6) \\
		r^3 & \mapsto (1\ 4\ 3\ 2)(5\ 6\ 7\ 8) & sr^3 & \mapsto (1\ 8)(2\ 5)(3\ 6)(4\ 7)
	\end{align*}
	\item The images of $D_8$ under the left regular representation with respect to this new labelling are as follows:
	\begin{align*}
		1 & \mapsto 1 & s & \mapsto (1\ 2)(3\ 4)(5\ 6)(7\ 8) \\
		r & \mapsto (1\ 3\ 5\ 7)(2\ 8\ 6\ 4) & sr & \mapsto (1\ 4)(2\ 7)(3\ 6)(5\ 8) \\
		r^2 & \mapsto (1\ 5)(3\ 7)(2\ 6)(4\ 8) & sr^2 & \mapsto (1\ 6)(2\ 5)(3\ 8)(4\ 7) \\
		r^3 & \mapsto (1\ 7\ 5\ 3)(2\ 4\ 6\ 8) & sr^3 & \mapsto (1\ 8)(2\ 3)(4\ 5)(6\ 7)
	\end{align*}
	Clearly, the two subgroups of $S_8$ obtained in parts (a) and (b) are different since, for example, the image of $r$ in part (a) is a product of two 4-cycles while the image of $r$ in part (b) is a product of two 4-cycles but with different elements.
\end{solenum}

\begin{exercise}[4.2.4]
	Use the left regular representation of $Q_8$ to produce two elements of $S_8$ which generate a subgroup of $S_8$ isomorphic to the quaternion group $Q_8$.
\end{exercise}

\begin{sol}
	At minimum, we have that $Q_8 = \gen{i, j}$ (any 2 elements of order 4 can be used), so we compute the images of $i$ and $j$ under the left regular representation.
	We label $Q_8$ as $1, -1, i, -i, j, -j, k, -k$ and label these with the integers $1,2 \ldots 8$ respectively.
	We find that
	\begin{align*}
		i & \mapsto \sigma_i = (1\ 3\ 4\ 2)(5\ 7)(6\ 8) \\
		j & \mapsto \sigma_j = (1\ 5\ 2\ 6)(3\ 8)(4\ 7)
	\end{align*}
	Hence $Q_8 \cong \gen{\sigma_i, \sigma_j} \leq S_8$.
\end{sol}

\begin{exercise}[4.2.5]
	Let $r$ and $s$ be the usual generators for the dihedral group of order 8 and let $H = \langle s \rangle$.
	List the left cosets of $H$ in $D_8$ as $1H, rH, r^2H$ and $r^3H$.
	\begin{subexercises}
		\item Label these cosets with the integers $1,2,3,4$, respectively.
		Exhibit the image of each element of $D_8$ under the representation $\pi_H$ of $D_8$ into $S_4$ obtained from the action of $D_8$ by left multiplication on the set of 4 left cosets of $H$ in $D_8$.
		Deduce that this representation is faithful (i.e., the elements of $S_4$ obtained form a subgroup isomorphic to $D_8$).
		\item Repeat part (a) with the list of cosets relabelled by the integers $1,3,2,4$, respectively.
		Show that the permutations obtained from this labelling form a subgroup of $S_4$ that is different from the subgroup obtained in part (a).
		\item Let $K = \langle sr \rangle$, list the cosets of $K$ in $D_8$ as $1K, rK, r^2K$ and $r^3K$, and label these with the integers $1,2,3,4$.
		Prove that, with respect to this labelling, the image of $D_8$ under the representation $\pi_K$ obtained from left multiplication on the cosets of $K$ is the same \emph{subgroup} of $S_4$ as in part (a) (even though the subgroups $H$ and $K$ are different and some of the elements of $D_8$ map to different permutations under the two homomorphisms).
	\end{subexercises}
\end{exercise}

\begin{solenum}
	\item We have $\pi_H(r) = (1\ 2\ 3\ 4)$ and $\pi_H(s) = (2\ 4)$.
	We continue in this way to find that the images of the elements of $D_8$ under $\pi_H$ are:
	\begin{align*}
		\pi_H(1) & = 1 & \pi_H(s) & = (2\ 4) \\
		\pi_H(r) & = (1\ 2\ 3\ 4) & \pi_H(sr) & = (1\ 2)(3\ 4) \\
		\pi_H(r^2) & = (1\ 3)(2\ 4) & \pi_H(sr^2) & = (1\ 3) \\
		\pi_H(r^3) & = (1\ 4\ 3\ 2) & \pi_H(sr^3) & = (1\ 4)(2\ 3)
	\end{align*}
	so that the subgroup $\gen{\pi_H(r), \pi_H(s)} \leq S_4$ is isomorphic to $D_8$.
	Since the kernel of $\pi_H$ is trivial, then the representation is faithful.
	\item The new labelling affords the permutations $\pi_H(r) = (1\ 3\ 2\ 4)$ and $\pi_H(s) = (3\ 4)$.
	We have the following images:
	\begin{align*}
		\pi_H(1) & = 1 & \pi_H(s) & = (3\ 4) \\
		\pi_H(r) & = (1\ 3\ 2\ 4) & \pi_H(sr) & = (1\ 3)(2\ 4) \\
		\pi_H(r^2) & = (1\ 2)(3\ 4) & \pi_H(sr^2) & = (1\ 2) \\
		\pi_H(r^3) & = (1\ 4\ 2\ 3) & \pi_H(sr^3) & = (1\ 4)(2\ 3)
	\end{align*}
	Moreover, the subgroup is different from that in part (a) since, for example, the image of $r$ in part (a) is $(1\ 2\ 3\ 4)$ while the image of $r$ in this part is $(1\ 3\ 2\ 4)$.
	\item Under $\pi_K$, we have $\pi_K(r) = (1\ 2\ 3\ 4)$ and $\pi_K(s) = (1\ 2)(3\ 4)$.
	With respect to this labeling, we have the subgroup $\wh K = \gen{(1\ 2\ 3\ 4), (1\ 2)(3\ 4)}$.
	Let $\wh H$ be the subgroup obtained in part (a).
	Note that $(1\ 2\ 3\ 4)$ is contained in both $\wh H$ and $\wh K$.
	Moreover, we have following:
	\[
		(1\ 2\ 3\ 4)(2\ 4)(1\ 4\ 3\ 2) = (1\ 2)(3\ 4) \in \wh H \quad \text{and} \quad (1\ 2\ 3\ 4)^2(1\ 2)(3\ 4) = (2\ 4) \in \wh K
	\]
	so that both subgroups contain the other generator of the other subgroup.
	Therefore, $\wh H = \wh K$.
\end{solenum}

\begin{exercise}[4.2.6]
	Let $r$ and $s$ be the usual generators for the dihedral group of order 8 and let $N = \langle r^2 \rangle$.
	List the left cosets of $N$ in $D_8$ as $1N, rN, sN$ and $srN$.
	Label these cosets with the integers $1,2,3,4$ respectively.
	Exhibit the image of each element of $D_8$ under the representation $\pi_N$ of $D_8$ into $S_4$ obtained from the action of $D_8$ by left multiplication on the set of 4 left cosets of $N$ in $D_8$.
	Deduce that this representation is not faithful and prove that $\pi_N(D_8)$ is isomorphic to the Klein 4-group.
\end{exercise}

\begin{sol}
	We have $\pi_N(r) = (1\ 2)(3\ 4)$ and $\pi_N(s) = (1\ 3)(2\ 4)$.
	Continuing in this way, we find that the images of the elements of $D_8$ under $\pi_N$ are:
	\begin{align*}
		\pi_N(1) & = 1 & \pi_N(s) & = (1\ 3)(2\ 4) \\
		\pi_N(r) & = (1\ 2)(3\ 4) & \pi_N(sr) & = (1\ 4)(2\ 3) \\
		\pi_N(r^2) & = 1 & \pi_N(sr^2) & = (1\ 3)(2\ 4) \\
		\pi_N(r^3) & = (1\ 2)(3\ 4) & \pi_N(sr^3) & = (1\ 4)(2\ 3)
	\end{align*}
	Since $\ker(\pi_N)$ contains the non-trivial element $r^2$, then the representation is not faithful.
	Moreover, we have $\pi_N(D_8) = \set{1, (1\ 2)(3\ 4), (1\ 3)(2\ 4), (1\ 4)(2\ 3)}$, which is isomorphic to the Klein 4-group.
\end{sol}

\begin{exercise}[4.2.7]
	Let $Q_8$ be the quaternion group of order 8.
	\begin{subexercises}
		\item Prove that $Q_8$ is isomorphic to a subgroup of $S_8$.
		\item Prove that $Q_8$ is not isomorphic to a subgroup of $S_n$ for any $n \le 7$. [If $Q_8$ acts on any set $A$ of order $\le 7$ show that the stabilizer of any point $a \in A$ must contain the subgroup $\langle -1 \rangle$.]
	\end{subexercises}
\end{exercise}

\begin{solenum}
	\item This is done in \exref{4.2.4}.
	\item Let $Q_8$ act on a set $A$ of order $n \leq 7$.
	For any $a \in A$, consider the orbit $\oo_a$ of $a$ under this action.
	By Proposition 4.2, we have $|\oo_a| = |Q_8 : (Q_8)_a|$, where $(Q_8)_a$ is the stabilizer of $a$ in $Q_8$.
	Since $|\oo_a|$ divides $|A| \leq 7$, then $|\oo_a|$ must be 1, 2, 4, or 7.
	However, since $Q_8$ has no subgroup of index 7, then $|\oo_a|$ cannot be 7.
	If $|\oo_a| = 4$, then $(Q_8)_a$ has order 2, but the only subgroup of order 2 in $Q_8$ is $\langle -1 \rangle$.
	If $|\oo_a| = 2$, then $(Q_8)_a$ has order 4, and the only subgroups of order 4 in $Q_8$ are $\langle i \rangle$, $\langle j \rangle$, and $\langle k \rangle$, all of which contain $\langle -1 \rangle$.
	If $|\oo_a| = 1$, then $(Q_8)_a = Q_8$, which also contains $\langle -1 \rangle$.
	Therefore, in all cases, the stabilizer $(Q_8)_a$ contains $\langle -1 \rangle$.
	Since this holds for all $a \in A$, then $\langle -1 \rangle$ is contained in the kernel of the action.
	The action is not faithful, hence $Q_8$ cannot be isomorphic to a subgroup of $S_n$ for any $n \leq 7$.
\end{solenum}

\begin{exercise}[4.2.8]
	Prove that if $H$ has finite index $n$ then there is a normal subgroup $K$ of $G$ with $K \le H$ and $|G : K| \le n!$.
\end{exercise}

\begin{sol}
	Let $G$ act by left multiplication on the set of left cosets of $H$ in $G$.
	Then we have the homomorphism $\phi : G \to S_n$, which is the permutation representation of $G$ on $G/H$.
	Let $K = \ker(\phi) \subseteq H$, since $g \in K$ if and only if $gH = H$.
	By the First Isomorphism Theorem, we have $G/K \cong \phi(G) \leq S_n$, so that $|G : K| = |G/K| = |\phi(G)|$ divides $n!$.
	Therefore, $|G : K| \leq n!$.
\end{sol}

\begin{exercise}[4.2.9]
	Prove that if $p$ is a prime and $G$ is a group of order $p^\alpha$ for some $\alpha \in \z^+$, then every subgroup of index $p$ is normal in $G$.
	Deduce that every group of order $p^2$ has a normal subgroup of order $p$.
\end{exercise}

\begin{sol}
	Let $H$ be a subgroup of $G$ with index $p$.
	Then $|G : H| = p$, so the action of $G$ on the left cosets of $H$ in $G$ gives a homomorphism $\phi : G \to S_p$.
	Since $|G| = p^\alpha$, then by Lagrange's Theorem, the order of $\phi(G)$ divides $p^\alpha$.
	However, the only subgroups of $S_p$ whose order divides $p^\alpha$ are the trivial group and groups of order $p$.
	Since $\phi(G)$ acts transitively on the $p$ cosets of $H$, then $\phi(G)$ cannot be trivial.
	Therefore, $|\phi(G)| = p$, which is prime, so $\phi(G)$ is cyclic and hence Abelian.
	The kernel of $\phi$ is a normal subgroup of $G$ contained in $H$.
	Since the image $\phi(G)$ has order $p$, then by the First Isomorphism Theorem, we have $|G : \ker(\phi)| = p$.
	But since $\ker(\phi) \subseteq H$ and both have index $p$, then $\ker(\phi) = H$.
	Therefore, $H$ is normal in $G$.

	To deduce that every group of order $p^2$ has a normal subgroup of order $p$, let $G$ be a group of order $p^2$.
	By Cauchy's Theorem, there exists an element of order $p$ in $G$, which generates a subgroup $H$ of order $p$.
	Since the index of $H$ in $G$ is $p$, by the previous result, $H$ is normal in $G$.
\end{sol}

\begin{exercise}[4.2.10]
	Prove that every non-Abelian group of order 6 has a nonnormal subgroup of order 2.
	Use this to classify groups of order 6. [Produce an injective homomorphism into $S_3$.]
\end{exercise}

\begin{sol}
	Note that $2 \mid 6$ and $3 \mid 6$.
	By Cauchy's Theorem, there exists subgroups $H$ and $K$ of orders 2 and 3 respectively.
	Let $H = \set{1, h}$, suppose it is normal, and let $g \in G - H$.
	Since $H \nsub G$, then $gHg\inv = H$.
	In particular, $ghg\inv \in H$, so either $ghg\inv = 1$ or $ghg\inv = h$.
	Since the first case implies $h = 1$, it must be that $ghg\inv = h$, or $gh = hg$.
	Then $h$ commutes with every element.
	Moreover, we may use \exref{3.3.3} to conclude that $G = HK$ since $H \nsub G$.
	Since $K$ is also Abelian as it is cyclic, then $G = hk$ for $h \in H$ and $k \in K$, implying that $G$ is Abelian, a contradiction.
	Therefore, $H$ is not normal in $G$.

	To classify groups of order 6, let $G$ be a group of order 6.
	If $G$ is Abelian, then $G \cong \z_6 \cong \z_2 \times \z_3$.
	If $G$ is non-Abelian, then by the above result, $G$ has a nonnormal subgroup $H$ of order 2.
	Let $K$ be the subgroup of order 3.
	Then $G = HK$.
	Let $G$ act by left multiplication on the left cosets of $K$ in $G$.
	This action gives a homomorphism $\phi : G \to S_3$.
	Since $H$ is not normal in $G$, then the action is non-trivial, so $\phi$ is injective.
	Therefore, $G \cong \phi(G) \leq S_3$.
	Since $|G| = 6 = |S_3|$, then $\phi(G) = S_3$, and hence $G \cong S_3$.
\end{sol}

\begin{exercise}[4.2.11]
	Let $G$ be a finite group and let $\pi : G \to S_G$ be the left regular representation.
	Prove that if $x$ is an element of $G$ of order $n$ and $|G| = mn$, then $\pi(x)$ is a product of $m$ $n$-cycles.
	Deduce that $\pi(x)$ is an odd permutation if and only if $|x|$ is even and $\abs G/\abs x$ is odd.
\end{exercise}

\begin{sol}
	Let $x \in G$ with $\abs x = n$.
	Consider the action of $\langle x \rangle$ on $G$ by left multiplication.
	The orbit of any $g \in G$ under this action is $\oo_g = \set{x^k g \mid k = 0, 1, \ldots, n-1}$.
	Since $\abs x = n$, then $|\oo_g| = n$ for all $g \in G$.
	By Proposition 4.2, we have $|\oo_g| = |\langle x \rangle : (\langle x \rangle)_g|$, where $(\langle x \rangle)_g$ is the stabilizer of $g$ in $\langle x \rangle$.
	Since $|\oo_g| = n$, then $(\langle x \rangle)_g$ is trivial for all $g \in G$.
	Therefore, the orbits of this action partition $G$ into subsets of size $n$.
	Since $|G| = mn$, there are exactly $m$ such orbits.
	Each orbit corresponds to an $n$-cycle in the permutation $\pi(x)$.
	Therefore, $\pi(x)$ is a product of $m$ $n$-cycles.

	To determine when $\pi(x)$ is an odd permutation, we note that an $n$-cycle is an odd permutation if and only if $n$ is even.
	Since $\pi(x)$ is a product of $m$ $n$-cycles, the parity of $\pi(x)$ is determined by the parity of $m$ and $n$.
	Specifically, $\pi(x)$ is odd if and only if $n$ is even and $m$ is odd.
	Thus, $\pi(x)$ is an odd permutation if and only if $|x|$ is even and $|G|/|x|$ is odd.
\end{sol}

\begin{exercise}[4.2.12]
	Let $G$ and $\pi$ be as in the preceding exercise.
	Prove that if $\pi(G)$ contains an odd permutation then $G$ has a subgroup of index 2. [Use \exref{3.3.3}.]
\end{exercise}

\begin{sol}
	Recall the $\ep$ homomorphism from Proposition 3.23.
	Moreover, $\pi$ is a homomorphism from $G$ to $S_G$.
	Consider the composition $\ep \circ \pi : G \to \set{\pm 1}$.
	Since $\pi(G)$ contains an odd permutation, then $\ep \circ \pi$ is surjective.
	By the First Isomorphism Theorem, we have $G/\ker(\ep \circ \pi) \cong \set{\pm 1}$, so that $|\ker(\ep \circ \pi)| = |G|/2$.
	Therefore, $\ker(\ep \circ \pi)$ is a subgroup of $G$ of index 2.
\end{sol}

\begin{exercise}[4.2.13]
	Prove that if $|G| = 2k$ where $k$ is odd then $G$ has a subgroup of index 2. [Use Cauchy's Theorem to produce an element of order 2 and then use the preceding two exercises.]
\end{exercise}

\begin{sol}
	By Cauchy's Theorem, there exists $x \in G$ such that $\abs x = 2$.
	Then $\pi(x)$ is of order 2 and is hence a product of disjoint transpositions.
	Since $|G| = 2k$ with $k$ odd, then use $m = k$ and $n = 2$ in \exref{4.2.11} along with even $\abs x$ to conclude that $\pi(x)$ is an odd permutation.
	By the \exref{4.2.12}, $G$ has a subgroup of index 2.
\end{sol}

\begin{exercise}[4.2.14]
	Let $G$ be a finite group of composite order $n$ with the property that $G$ has a subgroup of order $k$ for each positive integer $k$ dividing $n$.
	Prove that $G$ is not simple.
\end{exercise}

\begin{sol}
	Let $S$ be the set of all prime factors of $n$.
	By the Well Ordering Principle, this has a minimal element $p$.
	By hypothesis, $G$ has a subgroup $P$ of order $n/p$ with index $p$.
	By Corollary 4.5, $P$ is normal in $G$ since $p$ is the smallest prime dividing $n$.
	Therefore, $G$ is not simple.
\end{sol}